\documentclass[11pt]{article}

\usepackage[margin=1in]{geometry}
\usepackage{amsmath,amssymb,amsfonts}
\usepackage[hidelinks]{hyperref}

\newcommand{\Spin}{\mathrm{Spin}}
\newcommand{\SO}{\mathrm{SO}}
\newcommand{\SU}{\mathrm{SU}}
\newcommand{\U}{\mathrm{U}}
\newcommand{\CP}{\mathbb{CP}}
\newcommand{\cO}{\mathcal{O}}

\title{Route D1-F: F-theory Local Placement Audit}
\author{BoYu Chen}
\date{June 11, 2026}

\begin{document}
\maketitle

\section*{Purpose}

This note tests the first Route-D string-constrained graft: whether the
Route-A skeleton can be reproduced as a standard local F-theory GUT structure.
The target is deliberately local.  The note checks a codimension-one
\(D_5\) singularity for the \(\Spin(10)\) gauge algebra, a codimension-two
\(D_5\to E_6\) enhancement for a \(16\) matter curve, and a
\(\CP^1\) matter-curve index giving the same \(\cO(2)\) family carrier used in
the Route-A paper.  It does not construct a global compactification and does
not claim a unique string vacuum.

\section{Local F-theory placement}

In local F-theory GUT models, a gauge algebra is engineered on a complex
surface or divisor \(S\) by the singular fiber type over \(S\).  A \(D_5\)
singularity realizes the algebra \(\mathfrak{so}(10)\), whose simply connected
group is \(\Spin(10)\).  Charged matter is localized on curves where the
singularity enhances in codimension two.  The Katz--Vafa local rule determines
the localized matter by decomposing the adjoint of the enhanced algebra under
the lower-rank gauge algebra and its commutant.

For a \(D_5\to E_6\) enhancement, the relevant branching is
\begin{equation}
  E_6 \supset \SO(10)\times \U(1),
\end{equation}
and
\begin{equation}
  \mathbf{78}
  =
  \mathbf{45}_0
  \oplus
  \mathbf{1}_0
  \oplus
  \mathbf{16}_{-3}
  \oplus
  \overline{\mathbf{16}}_{+3}.
  \label{eq:e6-so10-branch}
\end{equation}
The dimension check is
\begin{equation}
  45+1+16+16=78.
\end{equation}
Thus the extra local degrees of freedom away from the
\(\mathfrak{so}(10)\oplus\mathfrak u(1)\) adjoint are a single \(16\) and its
conjugate.  In this sense the Route-A ``six-face'' object has a standard
string-local origin: the matter curve carries one irreducible \(\Spin(10)\)
half-spin representation, not six unrelated low-energy fragments.

\paragraph{Boundary.}
Equation~\eqref{eq:e6-so10-branch} is a local representation statement.  It
does not by itself determine net chirality, family number, Yukawa ranks, or
absence of exotics.  Those require flux data, global geometry, and a
compactification-specific spectrum computation.

\section{\texorpdfstring{Matter-curve index and the \(\cO(2)\) carrier}{Matter-curve index and the O(2) carrier}}

On a matter curve \(\Sigma\), the chiral zero-mode carrier in the local
twisted seven-brane theory is conventionally represented as a spin bundle on
\(\Sigma\) tensored with a flux line bundle:
\begin{equation}
  K_\Sigma^{1/2}\otimes L_{\rm flux}.
  \label{eq:spin-flux-carrier}
\end{equation}
For the Route-D test, choose
\begin{equation}
  \Sigma=\CP^1,\qquad
  K_{\CP^1}=\cO(-2),\qquad
  K_{\CP^1}^{1/2}=\cO(-1).
\end{equation}
If the restriction of the flux line bundle to the matter curve has degree
three,
\begin{equation}
  L_{\rm flux}|_\Sigma\simeq \cO(3),
\end{equation}
then
\begin{equation}
  K_{\CP^1}^{1/2}\otimes L_{\rm flux}
  =
  \cO(-1)\otimes\cO(3)
  =
  \cO(2).
  \label{eq:o2-from-flux}
\end{equation}
Therefore the Route-A family carrier is reproduced exactly, but with a
string-local decomposition of the degree:
\begin{equation}
  2=3_{\rm flux}-1_{\rm spin}.
\end{equation}
The cohomology on \(\CP^1\) is
\begin{equation}
  h^0(\CP^1,\cO(n))=
  \begin{cases}
    n+1,& n\ge0,\\
    0,& n<0,
  \end{cases}
  \qquad
  h^1(\CP^1,\cO(n))=
  \begin{cases}
    0,& n\ge -1,\\
    -n-1,& n\le -2.
  \end{cases}
\end{equation}
For \(n=2\),
\begin{equation}
  h^0(\CP^1,\cO(2))=3,\qquad
  h^1(\CP^1,\cO(2))=0.
\end{equation}
Thus the F-theory matter-curve audit recovers the same protected
three-section carrier used in Route A.

\paragraph{Boundary.}
The integer \(\deg L_{\rm flux}=3\) is still compactification data.  Route D
does not eliminate conditional input; it changes the input from an abstract
\(\cO(2)\) assumption into a flux-quantized matter-curve datum.

\section{Hypercharge flux and breaking caveat}

The face projection in the Route-A paper breaks the high-energy
\(\Spin(10)\) representation to the Standard-Model face.  In F-theory, this
kind of breaking can be implemented by internal fluxes or Wilson-line data,
but the hypercharge direction is constrained.  A viable hypercharge flux must
avoid giving the \(\U(1)_Y\) gauge boson a Stueckelberg mass; equivalently,
the corresponding flux must satisfy the relevant global topological
triviality condition after pushforward to the base.  It must also avoid
unwanted chiral exotics and keep threshold corrections under control.

Consequently, the D1-F audit supports the local origin of the
\(\Spin(10):16\) matter object and the \(\cO(2)\) carrier, but it does not
certify the full Standard-Model breaking sector.  Hypercharge flux is a
separate global audit, not an automatic consequence of the local
\(D_5\to E_6\) enhancement.

\section{Result}

The D1-F graft is viable in the limited Route-D sense:
\begin{align}
  D_5\text{ singularity on }S
  &\leadsto
  \Spin(10)\text{ gauge algebra},\\
  D_5\to E_6\text{ enhancement on }\Sigma
  &\leadsto
  \mathbf{16}\text{ matter curve},\\
  \Sigma=\CP^1,\quad
  K_\Sigma^{1/2}\otimes \cO(3)
  &=
  \cO(2),\\
  h^0(\CP^1,\cO(2))=3,\quad h^1(\CP^1,\cO(2))&=0.
\end{align}
This is strong enough for a conservative Route-D appendix paragraph, provided
the text keeps the global compactification and hypercharge-flux caveats.

\begin{thebibliography}{9}

\bibitem{KatzVafa1996}
S.~Katz and C.~Vafa,
``Matter from geometry,''
\href{https://arxiv.org/abs/hep-th/9606086}{arXiv:hep-th/9606086}.

\bibitem{BeasleyHeckmanVafa2008I}
C.~Beasley, J.~J.~Heckman, and C.~Vafa,
``GUTs and exceptional branes in F-theory I,''
\href{https://arxiv.org/abs/0802.3391}{arXiv:0802.3391}.

\bibitem{BeasleyHeckmanVafa2008II}
C.~Beasley, J.~J.~Heckman, and C.~Vafa,
``GUTs and exceptional branes in F-theory II: Experimental predictions,''
\href{https://arxiv.org/abs/0806.0102}{arXiv:0806.0102}.

\bibitem{MayrhoferPaltiWeigand2013}
C.~Mayrhofer, E.~Palti, and T.~Weigand,
``Hypercharge flux in IIB and F-theory: Anomalies and gauge coupling
unification,''
\href{https://arxiv.org/abs/1303.3589}{arXiv:1303.3589}.

\end{thebibliography}

\end{document}
